%% Preamble additions for the chapters generated from stepss-docs.
%%
%% They are kept apart from stepss_doc.tex because they exist only to satisfy
%% pandoc's LaTeX writer: longtable/booktabs for its tables, calc for the
%% \real{} in its column widths, and \tightlist for its lists. Changing the
%% generator's output is the reason to change this file, and nothing else is.

% The guide has always drawn a coloured rule around every link, which is
% hyperref's default under pdflatex. That was tolerable at a few dozen
% cross-references; the generated chapters carry about two hundred, and a page
% of boxed phrases is hard to read. Colour the text instead, dark enough to
% print as near-black.
\hypersetup{
  colorlinks=true,
  linkcolor=[RGB]{26,54,93},
  citecolor=[RGB]{26,54,93},
  urlcolor=[RGB]{43,108,176},
  linktoc=all
}

\usepackage{longtable}
\usepackage{booktabs}
\usepackage{array}
\usepackage{calc}
\usepackage{graphicx}

% pandoc emits \tightlist for a list whose items carry no blank line between
% them, and defines it only inside its own template.
\providecommand{\tightlist}{%
  \setlength{\itemsep}{0pt}\setlength{\parskip}{0pt}}

% A note box for what the site writes as :::note, :::caution and :::tip. One
% environment rather than three, because the printed page distinguishes them by
% the heading the directive already carries.
\newmdenv[
  linewidth=0.6pt,
  linecolor=black!45,
  backgroundcolor=black!3,
  roundcorner=3pt,
  innertopmargin=6pt,
  innerbottommargin=6pt,
  innerleftmargin=8pt,
  innerrightmargin=8pt,
  skipabove=\topsep,
  skipbelow=\topsep
]{docsnotebox}

\newenvironment{docsnote}[1]
  {\begin{docsnotebox}\noindent\textbf{#1}\par\vspace{0.4ex}}
  {\end{docsnotebox}}

% A screenshot is a picture of a window and gains nothing from a rule around
% it, but a wide table set in \small still wants to breathe.
\setlength{\LTpre}{\medskipamount}
\setlength{\LTpost}{\medskipamount}

% The generated chapters quote data records and shell sessions at their natural
% width, which is wider than this text block. Let a long verbatim line run into
% the margin rather than silently past the page edge.
\usepackage{fvextra}
\fvset{breaklines=true,breakanywhere=true,fontsize=\small}

% ---------------------------------------------------------------- contents
% book.cls gives a section entry 2.3em for its number, which fits "9.9" and
% not "17.10". With 42 chapters, every chapter past the ninth that has ten
% sections ran its number into its title.
\makeatletter
\renewcommand*\l@section{\@dottedtocline{1}{1.5em}{3.0em}}
\renewcommand*\l@subsection{\@dottedtocline{2}{4.5em}{3.6em}}
\makeatother

% ------------------------------------------------------------- line breaks
% A URL is one unbreakable word to TeX and several of them are wider than the
% text block. xurl lets \url break at any character rather than only at the
% punctuation url.sty allows.
\usepackage{xurl}

% Typewriter text is not hyphenated by default, so a long identifier such as
% add_keyblock_resource sets as one box and runs into the margin. htt lets
% those break too.
\usepackage[htt]{hyphenat}

% Room for the paragraph breaker to work with before it gives up and
% overfills a line.
\setlength{\emergencystretch}{3em}
